% Atur variabel berikut sesuai dengan identitas Anda

% nama
\newcommand{\name}{KEINDRA DANTE HIDAYAT}
\newcommand{\authorname}{Hidayat, Keindra Dante}
\newcommand{\nickname}{Keindra}
\newcommand{\advisor}{Astria Nur Irfansyah, S.T., M.Eng., Ph.D.}
\newcommand{\coadvisor}{PEMBIMBING 2}
\newcommand{\examinerone}{PENGUJI 1}
\newcommand{\examinertwo}{PENGUJI 2}
\newcommand{\examinerthree}{PENGUJI 3}
\newcommand{\headofdepartment}{KEPALA DEPARTEMEN}

% identitas
\newcommand{\nrp}{5022231236}
\newcommand{\advisornip}{198103252010121002}
\newcommand{\coadvisornip}{000000000000000000}
\newcommand{\examineronenip}{000000000000000000}
\newcommand{\examinertwonip}{000000000000000000}
\newcommand{\examinerthreenip}{000000000000000000}
\newcommand{\headofdepartmentnip}{000000000000000000}

% judul
\newcommand{\tatitle}{PERBANDINGAN KINERJA ALGORITMA FXLMS DAN FXNLMS PADA SISTEM \emph{ACTIVE NOISE CANCELLATION} (ANC) \emph{HEADPHONE} BERBASIS FPGA}
\newcommand{\engtatitle}{\emph{PERFORMANCE COMPARISON OF FXLMS AND FXNLMS ALGORITHMS IN AN FPGA-BASED HEADPHONE ACTIVE NOISE CANCELLATION (ANC) SYSTEM}}

% tempat
\newcommand{\place}{Surabaya}

% jurusan
\newcommand{\studyprogram}{Teknik Elektro}
\newcommand{\engstudyprogram}{Electrical Engineering}

% fakultas
\newcommand{\faculty}{Teknologi Elektro dan Informatika Cerdas}
\newcommand{\engfaculty}{Intelligent Electrical and Informatics Technology}

% singkatan fakultas
\newcommand{\facultyshort}{FTEIC}
\newcommand{\engfacultyshort}{ELECTICS}

% departemen
\newcommand{\department}{Teknik Elektro}
\newcommand{\engdepartment}{Electrical Engineering}

% kode mata kuliah
\newcommand{\coursecode}{EE234799}
